\documentclass{article}
\PassOptionsToPackage{numbers,sort&compress}{natbib}
\usepackage[eandd]{neurips_2026}
\makeatletter
\renewcommand{\@noticestring}{Anonymous research draft for internal review.}
\makeatother
\usepackage[T1]{fontenc}
\usepackage{amsmath,amssymb,amsthm,booktabs,microtype}
\usepackage[hidelinks]{hyperref}
\newtheorem{proposition}{Proposition}

\begin{document}
\title{Best-First Search for Minimum-Vote Feature Screening}
\author{Anonymous Author(s)}
\maketitle

\begin{abstract}
A fitted classifier can screen proposed feature deletions before one candidate is refit. For each subset, we count the validation panels on which masking preserves predictive performance within tolerance, then take the lowest count across feature restorations. Best-first search reads the next restoration of the candidate with the highest remaining upper bound. We give a necessary and sufficient stopping condition and determine when the search uses the fewest queries for the fixed restoration order. It attained this minimum on 90 of 173 nonempty candidate lists. On the other 83, the first candidate won and every rival scored zero. Across all 184 lists from 16 classification tasks, it preserved every exhaustive nomination and used 2,382 candidate--restoration evaluations against 2,582 for depth-first pruning. Single-point masking was faster and achieved higher mean compression within tolerance on both original task panels.
\end{abstract}

\section{Introduction}
A classifier trained on two interchangeable features may use only one of them. Masking that feature can destroy the fitted classifier's predictions even though retraining on the other feature would recover the signal. This distinction matters when several deletions are proposed and only one candidate can be checked by retraining. A screening rule decides which deletion gets that check. Wrappers spend fits to compare subsets directly~\cite{kohavi1997wrappers}. Masking offers a cheaper measurement with a different target~\cite{covert2021removing,hooker2019roar}.

We study a screen that votes on each proposed deletion with some of the other features restored. A candidate's lowest vote across restorations determines its rank. One bad restoration can therefore eliminate a candidate, whereas a high score on an early restoration leaves its rank unresolved. We use this asymmetry to choose the next prediction. The search is an application of best-first minimax~\cite{stockman1979,plaat1995} to a table with few possible vote counts and a fixed tie order. Our analysis identifies the observations needed for this particular feature-selection decision and measures their cost against depth-first pruning and single-point screening.

\section{Minimum-vote nomination}
\label{sec:minimum-vote}
Let \(A\) be the starting feature set and \(S_1,\ldots,S_n\subsetneq A\) the proposed retained sets in tie order. On the Gain slate, smaller sums of feature ranks come first, with remaining ties broken by the sorted column indices.
Fit a classifier once and mask hidden columns at medians computed on fitting rows.
For each restoration \(E\) of eligible columns outside \(A\), compare the predictions using \(S_j\cup E\) with those using \(A\cup E\).
The family \(\mathcal E\) starts with no restored columns and ends with all eligible columns outside \(A\). Its interior masks give each tail column a distinct binary signature of equal weight, so some mask separates each ordered pair of tail columns~\cite{spencer1970separating}.

Let \(Q_{bm}(D)\) denote metric \(m\) on proposal panel \(b\) with visible columns \(D\). The four panels use one weighted sample and three stratified bootstrap reweightings of the same rows. Write \(R_{bmE}=Q_{bm}(A\cup E)\). Counting panels that meet every tolerance gives the minimum vote
\[
 c_E(S)=\sum_{b=1}^{4}\mathbf 1\!\left\{
   R_{bmE}-Q_{bm}(S\cup E)\leq \tau_m |R_{bmE}|+10^{-12}
   \ \text{for all }m\right\},\quad
 v(S)=\min_{E\in\mathcal E}c_E(S).
\]
We use \(\tau_m=0.03\) for AUC and accuracy, and add average precision with tolerance \(0.05\) on binary tasks when the minority class comprises at most 5\% of the proposal rows. The minimum adapts the worst-coalition criterion of MinPI~\cite{shi2026minpi} to a proposed deletion. Nomination chooses the first maximizer of \(v(S)\), even when its score is zero.

\paragraph{Nomination and refit.}
For each candidate keep its running upper bound \(u_j\), initially four, and its position in the restoration list. The empty restoration is first.
\begin{enumerate}\setlength{\itemsep}{0pt}
\item Take the earliest candidate \(j\) with the largest \(u_j\).
\item Return \(S_j\) if its row is exhausted or \(u_j=0\). Also return \(S_1\) if \(j=1\) and every rival has upper bound zero, including the case of no rivals.
\item Otherwise evaluate the next restoration \(E\) for \(S_j\), set \(u_j\gets\min(u_j,c_E(S_j))\), and return to step 1.
\end{enumerate}
Cache predictions by the visible feature set so different candidates can reuse the same masked reference. Refit the nominee and the anchor on the fitting and proposal rows. If the nominee meets every tolerance on separate guard rows, deliver it. A failed check or an empty candidate list returns \(A\).

\section{The observations needed to select a winner}
\label{sec:certificate}
View the votes as a table with one row per candidate and one column per restoration. A query reveals one count. In the following result, unread counts may independently take any integer from zero to the number of panels \(B\). This is a model of the information available to the search. A fitted classifier can impose further constraints on its predictions.

\begin{proposition}
After any queries, let \(u_i\) be the smallest count read in row \(i\), initially \(B\). Set \(\ell_i=u_i\) for a fully read row and \(\ell_i=0\) otherwise. The observations determine \(w\) as the first maximum precisely when
\[
  u_i<\ell_w\quad(i<w),\qquad u_i\le\ell_w\quad(i>w).
\]
With a fixed restoration order within each row, best-first nomination uses the fewest queries that can establish the winner on the realized table whenever \(w>1\) or some rival has a positive final score.
\end{proposition}

\begin{proof}
Every compatible row minimum lies between \(\ell_i\) and \(u_i\). Both endpoints can be attained independently across rows: fill unread entries with zero for the lower endpoint and \(B\) for the upper. Putting the winner at its lower endpoint and every rival at its upper endpoint gives the two inequalities, including the strict one for earlier rivals.

Write \(s=v(S_w)\). Under the stated condition any sufficient collection of prefixes must exhaust row \(w\): with an unread entry its lower endpoint is zero, which cannot defeat an earlier rival strictly or a positive rival weakly. Each earlier rival must be read until its upper bound first falls below \(s\), and each later rival until its bound first reaches or falls below \(s\). A prefix of length zero can suffice. The true winner's running bound always remains at least \(s\). Once a rival reaches its required cutoff it ranks behind the winner and the search never reads it again. On termination the inequalities hold, so the search reads exactly the required prefixes.
\end{proof}

The exception is observable in hindsight. With two rows of length \(H\), let the first contain only \(B\) and the second only zeros. Best-first reads all \(H\) entries of the first row. Reading just the first zero in the second row also proves the first candidate wins. If both rows instead contain only \(B\), the unique shortest proof reads the first row alone. For \(H>1\) no first query starts a shortest proof on both tables. These counts concern votes read in fixed order. A cached reference prediction can serve several votes, and computation time also depends on the learner.

\section{Empirical evaluation}
\label{sec:minimum-results}
Ten UCI classification tasks with 4--41 predictors supplied 30 outer folds.
Six wider tasks with 72--856 predictors supplied 16, including chronological
splits for Ozone and SECOM, participant-level splits for Parkinson and the
official Urban split. Development splits allocated 75\%, 12.5\% and 12.5\%
to fitting, nomination and the refit check, by distinct dates in chronological
order for Ozone and SECOM and by stratified rows elsewhere. We ranked split-used features by XGBoost gain
and took the shortest tested prefix whose masked validation scores met the
8\% relative allowance as the anchor. We kept its leading half as one
candidate and sampled up to seven other subsets of the same size, rounding
the retained width up. Eleven singleton-anchor folds offered no deletion and remained in all panel means.

We report the fraction of anchor columns removed by the delivered model,
counting zero when it breaches an outer tolerance or returns the anchor.
We denote this effective compression by \(U\) and report the breach rate
separately. We averaged folds within tasks and gave each task equal weight.
Paired 95\% intervals resample whole tasks 30,000 times with their fitted
models and folds held fixed.

\begin{table}[tb]
\centering\small
\caption{Nomination cost and outer performance on identical Gain candidate
sets. Random forests performed both masking and refitting. \(U\) and breach
are task-weighted percentages. Masks count distinct predictions on the
nomination sample. Seconds sum within-fold medians from three repetitions.
The Gain prefix and uniform choice need no masked prediction.}
\label{tab:minimum-cost-outcome}
\setlength{\tabcolsep}{4pt}
\begin{tabular}{@{}lrrrrrrrr@{}}
\toprule
&\multicolumn{4}{c}{Ten tasks, 30 folds}&\multicolumn{4}{c}{Six wider tasks, 16 folds}\\
\cmidrule(lr){2-5}\cmidrule(l){6-9}
Rule & Masks & Sec. & \(U\) & Breach & Masks & Sec. & \(U\) & Breach\\
\midrule
Gain prefix     &0   &--- &1.56&3.33 &0     &--- &2.72 &16.67\\
Uniform choice  &0   &--- &4.79&6.81 &0     &--- &8.21 &12.96\\
Single point     &126 &0.61&6.33&10.00&125 &0.65&13.77&16.67\\
Depth-first vote &189 &0.89&2.89&3.33 &366 &1.96&8.28 &16.67\\
Best-first vote  &170 &0.81&2.89&3.33 &300 &1.63&8.28 &16.67\\
\bottomrule
\end{tabular}
\end{table}

\paragraph{Nomination cost.}
Depth-first pruning visits candidates in tie order and abandons each row as soon
as its upper bound cannot beat the incumbent. We kept both orders fixed, fitted a 100-tree forest per fold,
and timed each rule with a fresh prediction cache on one pinned CPU.
Each timed rule ran first, second and third once per candidate set. Both searches returned the
exhaustive vote's nominee on the 46 original sets and on 138 sets from
three draws of local feature exchanges around ANOVA-ranked anchors on the same folds. Exhaustive voting used 760
and 1,334 masks on the two original panels. Best-first search cut the
depth-first count by a further 10\% and 18\%. Nomination time fell by 9\%
and 17\%, respectively. A single point was faster on both panels
(Table~\ref{tab:minimum-cost-outcome}).
Best-first used fewer masked predictions than depth-first on 41 of 184 lists and tied
on 143. Of the 173 nonempty lists, 90 met the prefix minimum in Section~\ref{sec:certificate}.
On the other 83 the first candidate won and every rival scored zero. Those lists took
718 queries against a shortest total of 602. The clocks cover nomination, excluding ranking, model fitting,
the guard and the final refit.

\paragraph{Choice of subset.}
The single-point rule maximizes the smallest normalized tolerance margin
across metrics on the original panel with no features restored.
Selecting this score across all candidates increased
\(U\) over the minimum vote by 3.44 percentage points on the ten tasks
(\([-4.67,15.00]\)) and 5.49 on the six wider tasks (\([0.00,11.05]\)).
On the first panel its breach rate rose by 6.67 points. The first comparison
was sensitive to ties. Choosing uniformly among the highest-vote candidates
gave \(5.42\%\) effective compression, 2.53 points above the fixed order.

To isolate the contribution of voting with a common tie rule, we used
XGBoost for masking and refitting and retained the better half of each
fold's candidate sets by Gain rank sum. We compared the highest point
score on that shortlist with the highest point score among its
maximum-vote candidates. Voting changed only the Seeds nomination on
the ten tasks, improving \(U\) by 1.33 points (\([0,4.00]\)) with the
same breach rate. It changed none of the sixteen wider-fold nominations
and increased their total mask count from 71 to 168. This retrospective
comparison reused the measured outcomes. On these shortlists the vote
rarely altered the decision that received the refit.


\section{Scope}
Exactness concerns the winner of the specified ballot. The four panels reuse proposal rows, so their votes do not provide four independent measurements of risk. The one-candidate budget counts the nominee sent to the guard. For outer evaluation we refit the returned subset and its reference on all development rows and compare them on held-out rows. A nominee that passes the guard may still breach an outer tolerance. The benchmark reuses 16 purposively selected sources and one forest seed per fold. A stable prediction advantage on new sources has yet to be established.

\bibliographystyle{plainnat}
\bibliography{references,revision,ensemble-sources}
\end{document}
